\documentclass[11pt]{article}
\usepackage{geometry}
\geometry{verbose}
\setlength{\parskip}{6pt}
\setlength{\parindent}{0pt}
\usepackage{color}
\usepackage{amsmath, amssymb}
\usepackage[authoryear]{natbib}
\usepackage[unicode=true,
 bookmarks=false,
 breaklinks=false,pdfborder={0 0 1},backref=section,colorlinks=false]
 {hyperref}

\makeatletter
\newcommand{\R}{\mathbb{R}}
\newcommand{\E}{\mathbb{E}}
\newcommand{\Cov}{\mathrm{Cov}}
\newcommand{\Var}{\mathrm{Var}}
\providecommand{\tabularnewline}{\\}
\usepackage{graphicx}
\usepackage{float}
\usepackage{url}
\usepackage{array}
\usepackage{enumitem}

\title{\Huge Assignment 3}
\author{\Large Chris Conlon}
\date{\Large Fall 2026 --- due at the start of Session 6}

\makeatother

\begin{document}

\title{Assignment 3}
\maketitle
\begin{center}
\begin{tabular*}{0.9\textwidth}{@{\extracolsep{\fill}}@{\extracolsep{\fill}}l@{\extracolsep{\fill}}l@{\extracolsep{\fill}}l}
Econometrics I & $\qquad$ & Professor Chris Conlon\tabularnewline
NYU Stern &  & Email: \href{mailto:cconlon@stern.nyu.edu}{cconlon@stern.nyu.edu}\tabularnewline
\end{tabular*}
\par\end{center}

\textbf{Instructions.} This assignment covers endogeneity, simultaneity, and
instrumental variables (Sessions 5--6). You may use \texttt{R} (with \texttt{fixest};
\texttt{feols(y \~{} x | p \~{} z)} estimates 2SLS) or Python (with \texttt{pyfixest}).
Turn in (1) a PDF presenting your results and showing your work, and (2) your code.
As always: derive with the operator rules, then simulate to check yourself.

\subsection*{1 Why you cannot regress quantity on price}

A grocery chain sets weekly prices for a snack product in 500 store-weeks. Demand and
pricing are:
\begin{align*}
q_{i} &= 8 - 2 \ln p_{i}^{*} + \xi_{i} && \text{(demand, in logs: elasticity } -2\text{)}\\
\ln p_{i}^{*} &= 1 + 0.5\, c_{i} + 0.3\, \xi_{i} && \text{(pricing: costs, and a response to demand)}
\end{align*}
where $\xi_{i} \sim \mathcal{N}(0, 0.5^2)$ is a demand shock the chain observes but
you do not, and $c_{i} \sim \mathcal{N}(0, 0.5^2)$ is a wholesale cost shifter.
\begin{enumerate}[label=(\alph*)]
\item Before simulating: use the operator rules to derive $\Cov(\ln p^{*}, \xi)$ and
the sign of the bias in the OLS regression of $q$ on $\ln p^{*}$. Will OLS make demand
look too elastic or too inelastic? Give the economic story in one sentence.
\item Simulate 500 observations. Estimate the elasticity by OLS and compare to your
answer in (a).
\item Now use the cost shifter $c$ as an instrument for $\ln p^{*}$
(\texttt{feols(q \~{} 1 | lnp \~{} c)}). Compare the estimate to $-2$.
\item State the two IV conditions (relevance, exclusion) in the language of this
problem, and say how you would defend each to a skeptical seminar audience.
One of them is checkable in data --- which one, and how?
\item A colleague proposes instrumenting price with \emph{the same product's price at
other stores in the chain}. Under what assumption is that valid, and what plausible
feature of grocery chains breaks it?
\end{enumerate}

\subsection*{2 The randomized experiment that wasn't enough}

\emph{Based on Jeziorski, Leng \& Seiler (2025), ``Causal Inference with Endogenous
Price Response'' (working paper; CEPR DP 20252) --- see the DAG from Session 5.}

An online marketplace randomly assigns a product-page upgrade ($T_{i} \in \{0,1\}$,
probability $0.5$) to 5{,}000 sellers. The upgrade shifts demand directly, but sellers
also \emph{reprice} in response to the upgrade and to demand conditions:
\begin{align*}
p_{i} &= 1 + 0.3\, T_{i} + 0.4\, \xi_{i} + 0.5\, c_{i}\\
q_{i} &= 2 - 2\, p_{i} + 1 \cdot T_{i} + \xi_{i}
\end{align*}
with $\xi_{i}, c_{i} \sim \mathcal{N}(0,1)$ independent; $c_{i}$ is a seller-level
cost shifter. The \textbf{direct effect} of the upgrade --- the demand shift holding
price constant --- is $1$.
\begin{enumerate}[label=(\alph*)]
\item Simulate the data. Regress $q$ on $T$ alone. What number do you get, and
\emph{of what is it an unbiased estimate}? (Hint: substitute the price equation into
the demand equation.) When would a manager want exactly this number?
\item Now ``control for price'': regress $q$ on $T$ and $p$. You should find an
estimate strictly between (a) and the direct effect of $1$. Explain, using the DAG
from lecture, why this regression is biased for the direct effect \emph{even though
$T$ is randomized}. Which arrow does the damage?
\item Run the diagnostic from lecture: regress $p$ on $T$. What do you find, and what
does it tell you?
\item Recover the direct effect by instrumenting price with the cost shifter:
\texttt{feols(q \~{} T | p \~{} c)}. Compare to $1$.
\item Managerial summary, three sentences max: your product team says ``the upgrade
only moved sales by [your answer to (a)], so it's not worth the engineering cost.''
What do you tell them, and which of your four numbers belongs in which business
decision?
\item Re-run (a)--(d) with the repricing channel shut off ($p_{i} = 1 + 0.4\xi_{i} +
0.5c_{i}$). Which estimators now agree, and why does the diagnostic in (c) predict
this?
\end{enumerate}

\subsection*{3 The Wald estimator, by hand}

Consider the just-identified model $y_{i} = \beta_{0} + \beta_{1} x_{i} +
\varepsilon_{i}$ with a binary instrument $z_{i}$ satisfying $\Cov(z, \varepsilon) = 0$
and $\Cov(z, x) \neq 0$.
\begin{enumerate}[label=(\alph*)]
\item Starting from the model, take the covariance of both sides with $z$ (operator
rules!) and solve for $\beta_{1}$. You have just derived
$\beta_{1} = \Cov(z, y)/\Cov(z, x)$.
\item Show that this equals the \emph{Wald estimator}
$\left(\E[y|z=1] - \E[y|z=0]\right) / \left(\E[x|z=1] - \E[x|z=0]\right)$.
\item In one sentence: why does dividing the ``reduced form'' by the ``first stage''
undo the attenuation that a weak first stage would otherwise cause --- and what is the
price you pay when the first stage is very weak?
\end{enumerate}

\end{document}
